% Fuente única de clase: Recursos/2025-11-25.pdf, páginas 1--10.
En la clase 26 relacionamos el grado separable con las raíces de los polinomios mínimos y probamos el teorema del elemento primitivo. Después introducimos los campos fijos y la correspondencia de Galois, que invierte las inclusiones entre campos y subgrupos.

%%%%%%%%%%%%%%%%%%%%%%%%%%%%%%%%%%%%%%%%%%%%%%%%%%%%%%%%%%%%%%%%%%%%%%%%%%%%%%%
\section{Grado separable y generación por elementos separables}
%%%%%%%%%%%%%%%%%%%%%%%%%%%%%%%%%%%%%%%%%%%%%%%%%%%%%%%%%%%%%%%%%%%%%%%%%%%%%%%
%%%%%%%%%%%%%%%%%%%%%%%%%%%%%%%%%%%%%%%%%%%%%%%%%%%%%%%%%%%%%%%%%%%%%%%%%%%%%%%
\subsection{Raíces e inmersiones de una extensión simple}
%%%%%%%%%%%%%%%%%%%%%%%%%%%%%%%%%%%%%%%%%%%%%%%%%%%%%%%%%%%%%%%%%%%%%%%%%%%%%%%
% P01--P02.
Fijemos una clausura algebraica $\overline k$ de $k$. La definición de la clase 25 se extiende a cualquier extensión algebraica $F/k$ mediante
\[
[F:k]_s=\card\Hom_k(F,\overline k).
\]
Ahora esta cardinalidad puede ser infinita. El conjunto no es vacío por el \cref{lema:C24-inmersion}, y cambiar de clausura produce una biyección por composición con un isomorfismo sobre $k$.

\begin{lema}\label{lema:C26-simple-separable}
Sea $k(\alpha)/k$ una extensión algebraica simple, con polinomio mínimo $p$. Entonces $[k(\alpha):k]_s$ es el número de raíces distintas de $p$ en $\overline k$. En particular,
\[
[k(\alpha):k]_s\leq[k(\alpha):k],
\]
con igualdad si y sólo si $\alpha$ es separable sobre $k$.
\end{lema}
\begin{proof}
Sean $\alpha_1,\ldots,\alpha_s$ las raíces distintas de $p$. Consideremos
\[
\begin{aligned}
\Phi:\Hom_k(k(\alpha),\overline k)&\longrightarrow\{\alpha_1,\ldots,\alpha_s\},\\
i&\longmapsto i(\alpha).
\end{aligned}
\]
Está bien definida porque $p(i(\alpha))=i(p(\alpha))=0$: $i$ fija todos los coeficientes de $p$. Todo elemento de $k(\alpha)$ es una expresión polinómica en $\alpha$ con coeficientes de $k$, de modo que dos homomorfismos con la misma imagen de $\alpha$ coinciden en todo el campo. Así $\Phi$ es inyectiva.

Para cada raíz $\beta$ de $p$, el isomorfismo de extensiones simples de la clase 22 envía $\alpha$ a $\beta$ y fija a $k$. Su composición con la inclusión en $\overline k$ es una preimagen de $\beta$ bajo $\Phi$. Los diagramas de evaluación y de transporte son
\[
\begin{tikzcd}[column sep=large,row sep=large]
k[t]\arrow[r,"\epsilon_\alpha"]\arrow[d,two heads]&
k(\alpha)\arrow[r,"i"]&\overline k\\
k[t]/(p)\arrow[ur,"\sim"']&&
\end{tikzcd}
\qquad
\begin{tikzcd}[column sep=large,row sep=large]
&k(\alpha)\arrow[dd,"\sim"]\arrow[dr,"i",destacado]&\\
k\arrow[ur,hook]\arrow[dr,hook]&&\overline k\\
&k(\beta)\arrow[ur,hook]&
\end{tikzcd}
\]
La biyección da $[k(\alpha):k]_s=s\leq\deg p=[k(\alpha):k]$. La igualdad equivale a que todas las raíces de $p$ sean distintas.
\end{proof}

%%%%%%%%%%%%%%%%%%%%%%%%%%%%%%%%%%%%%%%%%%%%%%%%%%%%%%%%%%%%%%%%%%%%%%%%%%%%%%%
\subsection{La fórmula de la torre para el grado separable}
%%%%%%%%%%%%%%%%%%%%%%%%%%%%%%%%%%%%%%%%%%%%%%%%%%%%%%%%%%%%%%%%%%%%%%%%%%%%%%%
% P02--P03.
\begin{lema}\label{lema:C26-torre-separable}
Sean $k\subseteq E\subseteq F$ extensiones algebraicas. El grado $[F:k]_s$ es finito si y sólo si lo son $[F:E]_s$ y $[E:k]_s$. En ese caso,
\[
[F:k]_s=[F:E]_s[E:k]_s.
\]
\end{lema}
\begin{proof}
Incluyamos $F$ en $\overline k$. Este último campo es también una clausura algebraica de $E$: es algebraico sobre $E$ y algebraicamente cerrado. La restricción define
\[
R:\Hom_k(F,\overline k)\longrightarrow\Hom_k(E,\overline k),
\qquad i\longmapsto\restr{i}{E}.
\]
Es sobreyectiva. En efecto, si $j:E\to\overline k$ fija a $k$, el argumento de prolongación mediante Zorn del \cref{lema:C24-inmersion}, con el homomorfismo inicial $j$, lo extiende a un homomorfismo $F\to\overline k$. El esquema es
\[
\begin{tikzcd}[column sep=large,row sep=large]
F\arrow[dr,dashed,"i",destacado]&\\
E\arrow[u,hook]\arrow[r,"j"']&\overline k\\
k\arrow[u,hook]\arrow[ur,hook]&
\end{tikzcd}
\]
Todas las fibras de $R$ tienen cardinalidad $[F:E]_s$. Para verlo, prolongamos $j$ a un homomorfismo $\tau:\overline k\to\overline k$ que fija a $k$. Es un automorfismo: su imagen es algebraicamente cerrada, y $\overline k$ es algebraico sobre ella porque ya es algebraico sobre $k$. El \cref{lema:C24-cerrado} obliga a que la imagen sea todo $\overline k$. Por tanto
\[
\Hom_E(F,\overline k)\longrightarrow R^{-1}(j),
\qquad h\longmapsto\tau\circ h
\]
es una biyección, con inversa $i\mapsto\tau^{-1}\circ i$.

Si $\Hom_k(E,\overline k)$ es infinito, la sobreyectividad hace infinito al dominio de $R$. Si una fibra es infinita, también lo es el dominio. Si ambos grados del enunciado son finitos, contamos las fibras mediante la regla del producto y obtenemos la fórmula.
\end{proof}

\begin{prop}\label{prop:C26-separable-generadores}
Sea $F/k$ finita. Entonces $[F:k]_s\leq[F:k]$, y son equivalentes:
\begin{enumerate}[nosep]
\item $F=k(\alpha_1,\ldots,\alpha_r)$, donde cada $\alpha_i$ es separable sobre $k$;
\item $F/k$ es separable;
\item $[F:k]_s=[F:k]$.
\end{enumerate}
\end{prop}
\begin{proof}
Una base vectorial finita genera $F$ como campo. Para cualquier familia finita de generadores, pongamos $K_0=k$ y $K_i=k(\alpha_1,\ldots,\alpha_i)$. Las dos fórmulas de la torre y el lema para extensiones simples dan
\[
[F:k]_s=\prod_{i=1}^r[K_i:K_{i-1}]_s
\leq\prod_{i=1}^r[K_i:K_{i-1}]=[F:k].
\]
Si cada $\alpha_i$ es separable sobre $k$, también lo es sobre $K_{i-1}$. En efecto, su mínimo sobre $K_{i-1}$ divide al mínimo sobre $k$, que no tiene raíces múltiples. Todos los factores de la desigualdad son entonces iguales, lo que prueba $(1)\Rightarrow(3)$.

Supongamos (3) y tomemos $\alpha\in F$. En la torre $k\subseteq k(\alpha)\subseteq F$ tenemos
\[
[F:k]_s=[F:k(\alpha)]_s[k(\alpha):k]_s,
\qquad
[F:k]=[F:k(\alpha)][k(\alpha):k].
\]
Cada factor de la primera igualdad es positivo y menor o igual que el correspondiente de la segunda. La igualdad de los productos obliga a la igualdad de cada pareja. Así $[k(\alpha):k]_s=[k(\alpha):k]$, y el \cref{lema:C26-simple-separable} muestra que $\alpha$ es separable. Esto prueba $(3)\Rightarrow(2)$. Por último, si la extensión es separable, cualquier familia finita de generadores satisface (1).
\end{proof}

%%%%%%%%%%%%%%%%%%%%%%%%%%%%%%%%%%%%%%%%%%%%%%%%%%%%%%%%%%%%%%%%%%%%%%%%%%%%%%%
\section{El teorema del elemento primitivo}
%%%%%%%%%%%%%%%%%%%%%%%%%%%%%%%%%%%%%%%%%%%%%%%%%%%%%%%%%%%%%%%%%%%%%%%%%%%%%%%
%%%%%%%%%%%%%%%%%%%%%%%%%%%%%%%%%%%%%%%%%%%%%%%%%%%%%%%%%%%%%%%%%%%%%%%%%%%%%%%
\subsection{Finitud de los campos intermedios}
%%%%%%%%%%%%%%%%%%%%%%%%%%%%%%%%%%%%%%%%%%%%%%%%%%%%%%%%%%%%%%%%%%%%%%%%%%%%%%%
% P04--P05: el caso de base finita se propone revisar.
\begin{ejer}[Extensiones finitas de campos finitos]\label{ejer:C26-campo-finito-simple}
Probar que toda extensión finita de un campo finito es simple.
\end{ejer}
\begin{sol}
Sea $F/k$ finita, con $k$ finito. Una base muestra que $F$ es un conjunto finito. Probaremos la ciclicidad del grupo multiplicativo.

Sea $m$ el mínimo común múltiplo de los órdenes de sus elementos. Para cada primo $\ell$ que divide a $m$, elijamos un elemento cuyo orden tenga la máxima potencia $\ell^{a_\ell}$ presente en $m$. Elevándolo a la parte de su orden coprima con $\ell$, obtenemos un elemento $v_\ell$ de orden $\ell^{a_\ell}$. Los $v_\ell$ conmutan y sus órdenes son coprimos dos a dos; su producto $v$ tiene orden $m$ por la \cref{prop:C2-producto-orden}.

Todo elemento de $F^\times$ satisface $x^m=1$. Como un polinomio no nulo de grado $m$ tiene a lo sumo $m$ raíces, $|F^\times|\leq m$. La inclusión $\ang{v}\subseteq F^\times$ da la desigualdad opuesta, luego $F^\times=\ang{v}$. El campo $k(v)$ contiene todas las potencias de $v$ y a cero, por lo que coincide con $F$.
\end{sol}

\begin{prop}\label{prop:C26-intermedios-simple}
Una extensión algebraica $F/k$ es simple si y sólo si tiene un número finito de campos intermedios $k\subseteq E\subseteq F$.
\end{prop}
\begin{proof}
Supongamos $F=k(\alpha)$ y sea $q_k$ el mínimo de $\alpha$ sobre $k$. Para un campo intermedio $E$, sea $q_E$ el mínimo de $\alpha$ sobre $E$ y sea $E'$ el campo generado sobre $k$ por los coeficientes de $q_E$. Entonces
\[
k\subseteq E'\subseteq E\subseteq F,
\qquad q_E\in E'[t],\qquad q_E\mid q_k\quad\text{en }E[t].
\]
El polinomio $q_E$ sigue siendo irreducible sobre $E'$, pues una factorización allí también lo sería sobre $E$. Como $E'(\alpha)=E(\alpha)=F$, resulta
\[
[F:E']=[F:E]=\deg q_E.
\]
La torre da $[E:E']=1$, luego $E=E'$. Por tanto $q_E$ determina a $E$. Todos los $q_E$ son divisores mónicos de $q_k$ en el anillo $F[t]$, que es un DFU; un polinomio fijo tiene sólo finitos divisores mónicos. Hay, pues, finitos campos intermedios. Las relaciones utilizadas se resumen en
\[
\begin{tikzcd}[column sep=large,row sep=large]
F\arrow[r,equal]&k(\alpha)\\
E\arrow[u,hook]\arrow[r,hook]&E(\alpha)\arrow[u,equal]\\
k\arrow[u,hook]\arrow[r,hook]&k(\alpha)\arrow[u,equal]
\end{tikzcd}
\]

Recíprocamente, supongamos que hay finitos campos intermedios. Entre los subcampos de $F$ finitamente generados sobre $k$, elijamos uno maximal $K$, lo que es posible por esa finitud. Para todo $\beta\in F$, el campo $K(\beta)$ sigue siendo finitamente generado y contiene a $K$, de modo que $K(\beta)=K$. Así $F=K$ es finitamente generado y, por ser algebraico, es finito sobre $k$.

Si $k$ es finito, aplicamos el ejercicio anterior. Si $k$ es infinito, basta reducir una familia finita de generadores a uno, de dos en dos. Para $k(\alpha,\beta)$, los infinitos elementos $c\in k$ producen sólo finitos campos $k(c\alpha+\beta)$. Existen $c\ne c'$ tales que
\[
E=k(c\alpha+\beta)=k(c'\alpha+\beta).
\]
Entonces
\[
\alpha=\frac{(c'\alpha+\beta)-(c\alpha+\beta)}{c'-c}\in E,
\qquad \beta=(c\alpha+\beta)-c\alpha\in E.
\]
Por tanto $k(\alpha,\beta)=k(c\alpha+\beta)$. La misma propiedad de finitud vale dentro de cada campo generado por parte de la familia original. Por inducción, si $k(\alpha_1,\ldots,\alpha_{r-1})=k(u)$, podemos escribir $F=k(u,\alpha_r)=k(cu+\alpha_r)$ para algún $c\in k$.
\end{proof}

%%%%%%%%%%%%%%%%%%%%%%%%%%%%%%%%%%%%%%%%%%%%%%%%%%%%%%%%%%%%%%%%%%%%%%%%%%%%%%%
\subsection{Extensiones finitas separables y automorfismos}
%%%%%%%%%%%%%%%%%%%%%%%%%%%%%%%%%%%%%%%%%%%%%%%%%%%%%%%%%%%%%%%%%%%%%%%%%%%%%%%
% P06--P07.
\begin{prop}[Elemento primitivo]\label{prop:C26-primitivo}
Toda extensión finita separable $F/k$ es simple.
\end{prop}
\begin{proof}
Si $k$ es finito, el \cref{ejer:C26-campo-finito-simple} da el resultado. Supongamos que $k$ es infinito. Como $F$ tiene finitos generadores, basta probar que $F=k(\alpha,\beta)$ es simple y después proceder por inducción. Los campos intermedios son separables sobre $k$ por estar contenidos en $F$.

Sea $\mathcal I=\Hom_k(F,\overline k)$. Es finito, de cardinalidad $[F:k]_s$. Para $i\ne j$ en $\mathcal I$, los polinomios lineales
\[
i(\alpha)t+i(\beta),\qquad j(\alpha)t+j(\beta)
\]
son distintos: si coincidieran, $i$ y $j$ coincidirían en los generadores y en $k$, luego en todo $F$. Por ello
\[
P(t)=\prod_{\substack{i,j\in\mathcal I\\i\ne j}}
\bigl((i(\alpha)-j(\alpha))t+i(\beta)-j(\beta)\bigr)
\in\overline k[t]
\]
es no nulo. Si no hay parejas distintas, el producto es uno. El polinomio tiene finitas raíces y $k$ es infinito, así que podemos elegir $c\in k$ con $P(c)\ne0$.

Para $\gamma=c\alpha+\beta$, las imágenes $i(\gamma)$ son distintas cuando $i$ varía en $\mathcal I$. Todas son raíces del mínimo de $\gamma$ sobre $k$. En consecuencia,
\[
[F:k]_s\leq[k(\gamma):k]\leq[F:k].
\]
La separabilidad da $[F:k]_s=[F:k]$, por la \cref{prop:C26-separable-generadores}. Así $k(\gamma)$ y $F$ tienen la misma dimensión finita sobre $k$; la inclusión entre ellos obliga a que sean iguales.
\end{proof}

\begin{coro}\label{coro:C26-aut-normal}
Si $F/k$ es finita y separable, entonces
\[
|\Aut_k(F)|\leq[F:k],
\]
con igualdad si y sólo si $F/k$ es normal.
\end{coro}
\begin{proof}
Por el elemento primitivo, $F=k(\alpha)$. El \cref{coro:C22-automorfismos} identifica sus automorfismos con las raíces del mínimo $p$ de $\alpha$ que pertenecen a $F$. De ahí la desigualdad. Como $p$ es separable, la igualdad significa que todas sus raíces pertenecen a $F$; entonces $F$ es su campo de descomposición y es normal por el \cref{teo:C24-normal}. A la inversa, la normalidad hace que $p$ se descomponga en $F$, y la separabilidad da exactamente $\deg p=[F:k]$ raíces distintas.
\end{proof}

%%%%%%%%%%%%%%%%%%%%%%%%%%%%%%%%%%%%%%%%%%%%%%%%%%%%%%%%%%%%%%%%%%%%%%%%%%%%%%%
\section{Campos fijos y correspondencia de Galois}
%%%%%%%%%%%%%%%%%%%%%%%%%%%%%%%%%%%%%%%%%%%%%%%%%%%%%%%%%%%%%%%%%%%%%%%%%%%%%%%
%%%%%%%%%%%%%%%%%%%%%%%%%%%%%%%%%%%%%%%%%%%%%%%%%%%%%%%%%%%%%%%%%%%%%%%%%%%%%%%
\subsection{Las dos aplicaciones de la correspondencia}
%%%%%%%%%%%%%%%%%%%%%%%%%%%%%%%%%%%%%%%%%%%%%%%%%%%%%%%%%%%%%%%%%%%%%%%%%%%%%%%
% P07--P09.
\begin{defi}
Sea $F/k$ una extensión y $G\leq\Aut_k(F)$. El \emph{campo fijo} de $G$ es
\[
F^G=\set{\alpha\in F\given g(\alpha)=\alpha\text{ para todo }g\in G}.
\]
\end{defi}
Es un subcampo que contiene a $k$: los automorfismos fijan a $0$ y $1$, y preservan sumas, productos, opuestos e inversos. Así los elementos fijados simultáneamente por todos ellos son estables bajo esas operaciones.

\begin{defi}
La \emph{correspondencia de Galois} de $F/k$ consiste en las aplicaciones
\[
\begin{aligned}
\set{E\given k\subseteq E\subseteq F}&\longrightarrow
\set{G\given G\leq\Aut_k(F)},& E&\longmapsto\Aut_E(F),\\
\set{G\given G\leq\Aut_k(F)}&\longrightarrow
\set{E\given k\subseteq E\subseteq F},& G&\longmapsto F^G.
\end{aligned}
\]
\end{defi}
El grupo $\Aut_E(F)$ está formado por los automorfismos que fijan a $E$ punto por punto; es un subgrupo de $\Aut_k(F)$. Estas aplicaciones existen sin suponer finitud, normalidad ni separabilidad. Todavía no se afirma que sean inversas.

\begin{lema}\label{lema:C26-correspondencia}
La correspondencia de Galois invierte inclusiones:
\[
E_1\subseteq E_2\ \Longrightarrow\ \Aut_{E_2}(F)\subseteq\Aut_{E_1}(F),
\qquad
G_1\subseteq G_2\ \Longrightarrow\ F^{G_2}\subseteq F^{G_1}.
\]
Además,
\[
E\subseteq F^{\Aut_E(F)},\qquad G\subseteq\Aut_{F^G}(F).
\]
Si $E_1E_2$ es el subcampo de $F$ generado por $E_1$ y $E_2$, entonces
\[
\Aut_{E_1E_2}(F)=\Aut_{E_1}(F)\cap\Aut_{E_2}(F),
\qquad
F^{\ang{G_1,G_2}}=F^{G_1}\cap F^{G_2}.
\]
\end{lema}
\begin{proof}
Un automorfismo que fija a $E_2$ fija a cualquier subcampo $E_1\subseteq E_2$. Un elemento fijado por todo $G_2$ está fijado por todos los elementos de un subgrupo $G_1\subseteq G_2$. Esto prueba la inversión de inclusiones. Por definición, todos los automorfismos de $\Aut_E(F)$ fijan a cada elemento de $E$, y todos los elementos de $G$ fijan a $F^G$; resultan las otras dos inclusiones.

Un automorfismo fija a $E_1E_2$ si y sólo si fija a ambos campos: en la dirección no inmediata, todo elemento del campo generado es una expresión obtenida mediante sumas, productos y cocientes de elementos de $E_1$ y $E_2$. Por ello pertenece a ambos grupos de automorfismos exactamente cuando pertenece al de su compuesto.

Finalmente, si $x$ está fijado por $\ang{G_1,G_2}$, lo está por $G_1$ y $G_2$. Recíprocamente, si ambos subgrupos fijan a $x$, también lo hacen sus inversos y todos los productos de sus elementos. Como esas palabras generan $\ang{G_1,G_2}$, $x$ pertenece a su campo fijo.
\end{proof}

%%%%%%%%%%%%%%%%%%%%%%%%%%%%%%%%%%%%%%%%%%%%%%%%%%%%%%%%%%%%%%%%%%%%%%%%%%%%%%%
\subsection{Una correspondencia que no es biyectiva}
%%%%%%%%%%%%%%%%%%%%%%%%%%%%%%%%%%%%%%%%%%%%%%%%%%%%%%%%%%%%%%%%%%%%%%%%%%%%%%%
% P10.
\begin{ejem}\label{ejem:C26-cubica-real}
Sea $a=\sqrt[3]{2}$ la raíz real y $F=\mathbb Q(a)$. Eisenstein con $2$ da $[F:\mathbb Q]=3$. Por la fórmula de la torre, un campo intermedio tiene grado uno o tres sobre $\mathbb Q$; las únicas posibilidades son el campo base y el campo completo.

El campo $F$ está contenido en $\mathbb R$ y $t^3-2$ tiene allí una sola raíz, $a$. Todo automorfismo sobre $\mathbb Q$ debe fijarla; como genera $F$, el grupo $\Aut_{\mathbb Q}(F)$ es trivial. Los dos campos intermedios se envían al mismo subgrupo $\{\id\}$, mientras que el campo fijo de éste es $F$:
\[
\begin{array}{c|c}
E&\Aut_E(F)\\\hline
\mathbb Q&\{\id\}\\
F&\{\id\}
\end{array}
\qquad F^{\{\id\}}=F.
\]
La aplicación de campos a subgrupos no es inyectiva, y la de subgrupos a campos no alcanza a $\mathbb Q$. Esta extensión es separable, por ser de característica cero, pero no es normal.
\end{ejem}

\begin{obs}[Referencias de lectura]
El grado separable se estudia en \citet[cap.~VII, lemas 4.22--23 y proposición 4.24]{aluffi2009}. Para el elemento primitivo, los campos intermedios y el conteo de automorfismos, véanse VII.5.18--20. La comprobación para campos finitos se desarrolla en el ejercicio de esta clase. Los campos fijos y las propiedades iniciales de la correspondencia se encuentran en VII.6.1--4; el último ejemplo corresponde a VII.6.4.
\end{obs}
